\section{Conclusion and Recommendations}
\label{submission:sec:conclusion}

In this work, we have presented a rigorous methodological deconstruction and independent robustness audit of Quantization-Aware Model Merging (Q-Merge). Adopting a highly critical and constructive perspective, we systematically analyzed current evaluation practices, exposing crucial blind spots that undermine the deployment feasibility of test-time weight-space fusion under low-bit quantization. 

Our multi-axial empirical evaluations have revealed that continuous merging coefficients optimized under a simulated quantization operator overfit intensely to the exact mathematical rounding boundaries of that operator. Deploying these configurations across hardware-target boundaries triggers catastrophic representation collapse, often reducing performance to random-guess levels. Furthermore, we showed that the standard unsupervised prediction entropy minimization objective is highly fragile under realistic calibration shifts, collapsing under severe class skew, and that derivative-free black-box search (1+1 ES) outperforms gradient-based STE on the source optimization schema but exhibits significantly worse generalization and a larger generalization gap under schema mismatch.

To guide the community toward more robust and generalizable model-merging practices, we propose the following four methodological mandates:
\begin{enumerate}
    \item \textbf{Mandatory Cross-Operator Validation:} Researchers must evaluate learned merging coefficients under a variety of hardware-relevant deployment schemas (e.g., symmetric/asymmetric per-tensor representations) rather than restricting testing to the simulated source schema.
    \item \textbf{Calibration Stream Heterogeneity Audits:} Test-time adaptation and model merging methods must be stress-tested under non-idealized streaming characteristics, specifically auditing robustness under input noise and severe label skewness. Under severe class skew, standard unsupervised prediction entropy minimization collapses because it is blind to class labels and aggressively overfits to the dominant class representations. To resolve this structural vulnerability, we advocate for the research of \emph{confidence-thresholded pseudo-labeling} or \emph{self-supervised/contrastive adaptation objectives} that do not rely on raw prediction probability profiles. Pseudo-labeling with a strict confidence threshold (e.g., only updating parameters using samples where prediction confidence exceeds $0.95$) ensures that gradient updates are driven by high-certainty representations, preventing the optimizer from collapsing into degenerate, constant-class shortcut states. Alternatively, integrating self-supervised contrastive objectives (such as SimCLR or Barlow Twins proxies) can align representations without any task-label supervision, ensuring high robustness to class skew in the calibration stream.
    \item \textbf{Re-evaluation of Optimizer Paradigms:} The community should look beyond standard straight-through gradient approximations (STE) and explore derivative-free, black-box, or hybrid optimizer pipelines that can stably navigate the discontinuous and non-convex loss landscapes induced by low-bit rounding operators. Specifically, we advocate for the research of \emph{hybrid optimization pipelines}, which formally combine first-order gradient approximations with zero-order black-box optimization. Algorithmically, the optimizer begins with Straight-Through Estimation (STE) for $T_{\text{coarse}}$ steps to rapidly identify a low-loss basin, updating coefficients via: $\Lambda^{(t+1)} = \text{Proj}_{[0, 1]} \left( \Lambda^{(t)} - \eta \tilde{\nabla}_{\Lambda} \mathcal{L}(\Lambda^{(t)}) \right)$, where $\tilde{\nabla}$ is the straight-through gradient. Once the entropy loss plateau is reached (i.e., $| \mathcal{L}(\Lambda^{(t)}) - \mathcal{L}(\Lambda^{(t-1)}) | < \epsilon$), the pipeline switches to a local $(1+1)$-ES search for $T_{\text{fine}}$ steps, perturbing coefficients with adaptive Gaussian noise $z \sim \mathcal{N}(0, \sigma^2 I)$ and selecting candidates that minimize both entropy and the Total Variation (TV) spatial regularizer: $\mathcal{L}_{\text{total}}(\Lambda) = \mathcal{L}_{\text{entropy}}(\Lambda) + \alpha \mathcal{R}_{\text{TV}}(\Lambda)$. This hybrid approach marries the coarse convergence speed of STE with the precise local search and spatial smoothing of black-box strategies, locating generalization-resilient coefficient configurations that survive hardware target shifts. While our 1+1 ES comparator established a useful black-box baseline, future work should explore more complex, population-based black-box optimizers, such as Covariance Matrix Adaptation Evolution Strategy (CMA-ES) \cite{hansen2001cma} or Natural Evolution Strategies (NES) like OpenES \cite{salimans2017evolution}. Because CMA-ES dynamically models the covariance structure of candidate solutions, it is uniquely suited to capture parameter correlations across adjacent layer coefficients and navigate the rugged, narrow valleys of a quantized merging landscape. Crucially, by learning the shape of the local loss basin, CMA-ES guides candidate perturbation steps along directions that are aligned with the underlying continuous landscape, avoiding search directions that align too closely with a single operator's discrete rounding grids. Furthermore, population-based search maintains diversity, preventing premature convergence to a single simulated operator's brittle rounding boundary and potentially mitigating the cross-operator overfitting we have exposed. Additionally, future works should investigate how advanced, reconstruction-based post-training quantization (PTQ) rounding mechanisms (such as AdaRound \cite{nagel2020adaround} or BRECQ \cite{li2021brecq}) interact with learned merging coefficients to mitigate discretization noise.
    \item \textbf{Mitigating Task Interference Prior to Discretization:} Since extreme weight-space divergence between expert models induces high non-convexity in the multi-task loss landscape, future works should employ unquantized merging techniques that explicitly filter or resolve parameter conflicts (such as TIES-Merging or DARE) before imposing quantization constraints, thereby smoothing the landscape for more stable low-bit search.
\end{enumerate}

\textbf{Generalization to Billion-Parameter LLMs and VLMs:} While our empirical deconstruction was conducted on a lightweight model (\texttt{ViT-Tiny}, 5.7M parameters) to guarantee complete scientific control and eliminate confounding hyperparameter effects, our core insights have profound implications for modern billion-parameter architectures (e.g., Llama-3, CLIP, or InstructBLIP) as well as alternative architectural representations. Specifically, extending our deconstruction protocol to alternative lightweight architectural paradigms (such as convolutional-based MobileNetV3 or ResNet-18) would verify whether the observed cross-operator collapse is a universal phenomenon of weight-space discretization or is partially moderated by specific self-attention structural representations. In convolutional networks, highly localized, translation-invariant spatial kernels exhibit different parameter distributions and are often more sensitive to quantization-induced activation shifting due to smaller, fixed receptive fields. In contrast, Vision Transformers rely on global self-attention matrices with large dynamic ranges in projection layers. Here, quantization noise in query-key-value projections can propagate exponentially through the softmax attention maps, potentially making ViT representations structurally more fragile under low-bit rounding mismatches. Evaluating both paradigms side-by-side would provide crucial empirical answers regarding how receptive field structures alter the geometric properties of the quantized merging landscape. 

Additionally, model scale introduces a complex interplay between representations and discretization. In larger models (such as \texttt{ViT-Base}, \texttt{ViT-Large}, or multi-billion parameter LLMs), the increased parameter capacity often leads to smoother local loss basins in full precision due to over-parameterization, potentially easing weight-space search. However, this blessing is accompanied by a severe dimensionality curse in the quantized regime: as the parameter count scales, the number of independent discrete rounding thresholds increases exponentially, leading to an extremely dense and rugged boundary landscape. Under unconstrained test-time adaptation, the risk of overfitting to the specific simulated operator's scale and zero-point parameters is substantially magnified in larger models. This suggests that while full-precision ensembling may become more stable with scale, the Cross-Schema Generalization Gap is likely to expand when deploying larger, quantized merged models on heterogeneous edge hardware.

In large-scale systems, weights are frequently compressed using structured, low-bit group-wise or block-wise post-training quantization backends (such as AWQ \cite{lin2023awq} or GPTQ \cite{frantar2022gptq}). As the number of model parameters scales up, the dimensionality of the coefficient search space increases, which makes unregularized first-order gradient approximations like the STE even more prone to chaotic noise propagation across layer boundaries. Moreover, modern group-wise quantization introduces highly customized, channel-grouped scaling factors. Optimizing merging coefficients on simulated group-wise quantization will overfit intensely to the localized scaling factors. When deployed across diverse physical hardware accelerators (e.g., Apple Neural Engine, Edge TPUs, or NVIDIA TensorRT backends) that employ slightly different grouping configurations or hardware-specific rounding logic, the Cross-Schema Generalization Gap is expected to expand, potentially causing sudden, severe representation collapse in LLM conversational coherence or VLM alignment. Therefore, extending our proposed deconstruction protocols to alternative architectures and multi-billion parameter regimes remains a critical and urgent direction for the community.

Furthermore, while our audit focused primarily on weight-only quantization (W4), real-world edge execution on physical neural processing units (NPUs) and digital signal processors (DSPs) often requires joint weight-activation quantization (e.g., W4A8 or W4A4) to achieve optimal computational throughput and VRAM efficiency. In Vision Transformers, the self-attention mechanisms are highly susceptible to dynamic activation outliers, which introduce high-frequency discretization noise and signal clipping. This additional dynamic activation noise heavily distorts the geometric properties of the multi-task loss landscape, potentially expanding the Cross-Schema Generalization Gap. Under joint quantization, the continuous merging coefficients must navigate both the discrete rounding boundaries of the weights and the input-dependent clipping thresholds of the activations. Integrating outlier-aware activation quantization smoothers, such as SmoothQuant \cite{xiao2023smoothquant}, into the model-merging pipeline prior to weight discretization is a promising structural defense to migrate activation outliers to the more robust weight space, thereby stabilizing search and enhancing deployment-time generalizability.

Ultimately, we argue that the field of weight-space consolidation must transition away from over-optimistic "state-of-the-art" parameter chasing. By embracing rigorous, multi-axial, and deployment-realistic evaluation protocols, the deep learning community can build truly robust multi-task models that generalize from simulation to physical edge hardware.
